\chapter{Supplementary Spectroscopy Analysis and Diagnostics}
\label{app:spectroscopy-supplement}

This appendix is the audit trail for Chapter
\ref{chap:spectroscopy-application}.  It records the dataset identity, every
preprocessing and validation rule, the complete development and sensitivity
summaries, and the numerical diagnostics needed to distinguish finite
prediction from reliable Krylov coordinates.  The chapter and appendix are
derived from one frozen archive; the material below does not redefine the
primary analysis after observing the benchmark results.

\section{Data provenance and structural audit}
\label{sec:app-spectroscopy-provenance}

The source is the 2002 pharmaceutical-tablet NIR Shoot-out
\parencite{hopkins_2003,eigenvector_tablets}.  The archive was downloaded from
the Eigenvector Research data page on 8 September 2026.  Its identifiers are:
\begin{itemize}
    \item ZIP SHA-256:\par
    {\footnotesize\ttfamily
    04c03a4314ab8e8f63e0ecc788b657b5d1b5450791ce465ad41ad90cb5d52125};
    \item contained MAT-file SHA-256:\par
    {\footnotesize\ttfamily
    129a32ec9e194e568cbd96a200c88aabeb347156081b5dfbb0120372fdd6c22a}.
\end{itemize}
The distributed file contains 155 calibration, 40 validation, and 460 test
rows, giving 655 rather than the catalogue's stated 654 tablets.  All rows and
all spectral entries are finite, and there are no exact duplicate spectra
within either instrument.  Every complete row is retained.  Because the file
does not supply global tablet identifiers, immutable analysis identifiers join
the source split to one-based row order, for example
\texttt{calibrate:001}.  Corresponding instrument-1 and instrument-2 rows are
treated as paired observations.

Table~\ref{tab:app-spectroscopy-splits} quantifies the response imbalance that
is visible in Figure~\ref{fig:spectroscopy-overview}.  The validation response
has sample standard deviation 2.393 and range 188.2--199.5, whereas the
calibration and benchmark distributions are much wider.  In fact, 267 of the
460 benchmark assay values fall outside the validation range.  Validation-set
$R^2$ therefore has a small denominator and is not comparable with benchmark
$R^2$.

\begin{table}[htbp]
    \centering
    \small
    \caption[Assay by supplied archive split]{Assay response by supplied archive split.  Standard deviations
    use denominator $n-1$.}
    \label{tab:app-spectroscopy-splits}
    \begin{tabular}{@{}lrrrrr@{}}
        \toprule
        Split & $n$ & Mean & Standard deviation & Minimum & Maximum \\
        \midrule
        Calibration & 155 & 192.885 & 21.980 & 151.6 & 239.1 \\
        Validation  &  40 & 194.795 &  2.393 & 188.2 & 199.5 \\
        Benchmark   & 460 & 188.400 & 15.763 & 154.3 & 237.7 \\
        \bottomrule
    \end{tabular}
\end{table}

The wavelength units and domain are supplied by the external description; they
are not embedded as ordinary numeric metadata in the MAT object.  The archive
and catalogue provide no explicit open-data licence.  The frozen empirical
archive therefore contains checksums and a reproducible preparation workflow,
but does not redistribute the source MAT file.

\section{Functional representation and preprocessing}
\label{sec:app-spectroscopy-preprocessing}

Each spectrum is represented by its $T=650$ evaluations on the uniform
600--1898~nm grid.  The discrete geometry is
\begin{equation}
    \langle f,g\rangle_T=\frac{1}{T}\sum_{j=1}^{T} f_jg_j,
    \qquad T=650.
    \label{eq:app-spectroscopy-inner-product}
\end{equation}
For a training sample $\mathcal A$, wavelength and response means are computed
using only rows in $\mathcal A$.  The centred estimator has intercept
\begin{equation}
    \widehat\alpha_{\mathcal A}
      =\overline Y_{\mathcal A}
       -\frac{1}{T}\overline X_{\mathcal A}^{\top}
        \widehat\beta_{\mathcal A},
    \label{eq:app-spectroscopy-intercept}
\end{equation}
which is algebraically equivalent to the prediction form in
\eqref{eq:spectroscopy-prediction}.  The training means are applied unchanged
to held-out observations.  The source code also archives the resulting
original-input coefficient vector and intercept for every fitted model, so its
stored prediction can be reconstructed directly.

No primary operation other than this fold-local centring is applied.  In
particular, there is no variance scaling, smoothing, derivative construction,
wavelength selection, standard-normal-variate transformation, multiplicative-
scatter correction, or baseline correction.  These choices matter because
every such operation can add regularisation or alter the covariance geometry.

The primary final fit has $650>195$ and numerical covariance rank 194 after
centring.  At least 456 grid-coordinate null directions therefore remain, so a
unique unrestricted physical slope is not identified.  Coefficient vectors
are archived for replay and numerical comparison, not interpreted as recovered
chemical effects.

Three alternatives were fixed before benchmark fitting:
\begin{enumerate}
    \item instrument~1 restricted to 600--1798~nm inclusive, retaining 600
    wavelengths;
    \item instrument~1 smoothed row by row with a 15-point cubic
    Savitzky--Golay filter, derivative order zero and interpolation boundary
    mode; and
    \item the complete 650-wavelength instrument-2 spectrum for the same
    tablets.
\end{enumerate}
Each transformed dataset is centred within the relevant training sample.  The
fixed smoothing operator itself does not estimate parameters from held-out
rows; all sample-dependent centring remains fold local.

\section{Evaluation design and leakage safeguards}
\label{sec:app-spectroscopy-validation}

Table~\ref{tab:app-spectroscopy-validation-design} summarises the complete
evaluation design.  Random folds are formed separately within the supplied
calibration and validation sources.  Within each source, observations are
ordered by assay in groups of five and assigned one of the five fold labels at
random.  Every development outer fold thus contains exactly 31 calibration
and 8 validation tablets.  This preserves source composition while avoiding
folds dominated by one part of the response range.

\begin{table}[htbp]
    \centering
    \small
    \caption[Frozen spectroscopy evaluation design]{Frozen spectroscopy evaluation design.  Exact assignments and
    resampling matrices, rather than only seeds, are included in the archive.}
    \label{tab:app-spectroscopy-validation-design}
    \begin{tabular}{@{}p{3.0cm}p{10.9cm}@{}}
        \toprule
        Stage & Frozen rule \\
        \midrule
        Repeated development
          & Five repeats of five source-specific, assay-balanced outer folds;
            one complete out-of-fold metric per repeat. \\
        Inner selection
          & New five-fold balanced partitions inside each outer training set;
            Raw and Arnoldi share indices but choose independently. \\
        Order sensitivity
          & Five source-specific contiguous outer folds, each holding out 31
            calibration and 8 validation rows. \\
        Original split
          & Fit and tune on 155 calibration tablets; evaluate once on the 40
            validation tablets. \\
        Final development fit
          & Fit on all 195 development tablets; Raw and Arnoldi use one frozen
            common balanced inner partition. \\
        Designated benchmark
          & Predict the 460 supplied test tablets once; no benchmark outcome
            enters centring, tuning, repair, or variant selection. \\
        Conditional bootstrap
          & 10,000 paired iid draws and 10,000 circular moving-block draws at
            each block length 5, 10, and 20; models are held fixed. \\
        Randomness
          & Master seed 20260908 with separate deterministic streams for outer
            folds, inner folds, final tuning, iid bootstrap, and block
            bootstrap. \\
        \bottomrule
    \end{tabular}
\end{table}

Across five repeats, $195\times5=975$ out-of-fold predictions are stored for
each method.  These are repeated predictions of the same 195 tablets, not 975
independent observations.  Correspondingly, repeat summaries and fold
distributions are descriptive diagnostics rather than a basis for standard
independent-sample tests.  The development completion record additionally
states \texttt{benchmark\_variables\_loaded=false}: the development stage did
not load the test variables.  Only after its checks passed was the full stage
allowed to fit the 195-row models and evaluate the benchmark.

The benchmark bootstrap resamples tablet indices in common across methods.
This preserves pairing and estimates conditional variation over the supplied
benchmark empirical distribution.  It does not refit centring, tuning, or the
models and therefore excludes development-sample and selection uncertainty.
The circular block calculations preserve local row-order sequences, but no
batch or acquisition-order metadata are available to determine whether archive
order has a physical interpretation.

\section{Complete procedures and tuning rules}
\label{sec:app-spectroscopy-tuning}

The estimator formulas are given in Chapter~\ref{chap:methods}; Table
\ref{tab:app-spectroscopy-procedures} records the application-specific tuning
and diagnostics.  The requested path maximum is 70, further limited by grid
size, centred training-sample size, numerical rank, and recurrence breakdown.

\begin{table}[htbp]
    \centering
    \scriptsize
    \caption[Regularisation rules in the empirical application]{Estimator-specific regularisation rules in the frozen application.
    The procedures retain different finite-$m$ objectives and stopping rules.}
    \label{tab:app-spectroscopy-procedures}
    \begin{tabular}{@{}p{2.4cm}p{6.3cm}>{\raggedright\arraybackslash}p{5.1cm}@{}}
        \toprule
        Method & Regularisation choice & Principal diagnostics \\
        \midrule
        CG--FPLS
          & Released-code conjugate-residual recurrence and first discrepancy
            crossing; $\tau=1.01$, $\delta=0.10$, with an internal moment-GCV
            FPCR variance pilot.
          & Threshold, crossing status, moment-residual path, pilot dimension,
            and variance estimate. \\
        Raw FPLS
          & Minimum five-fold inner prediction MSE over the unscaled raw Krylov
            power sequence, fitted by response-space normal equations.
          & Candidate validity, raw-basis and response-design condition numbers,
            and Krylov column norms. \\
        Arnoldi FPLS
          & Independent minimum on the same inner folds, using twice
            reorthogonalised Gram--Schmidt and response-space QR.
          & Effective dimension, rank breakdown, response-design condition,
            and basis orthogonality defect. \\
        FPCR
          & Response-space GCV over covariance-eigenvalue cut-offs.
          & Numerical rank, retained eigenvalue, and retained spectral
            condition number. \\
        \bottomrule
    \end{tabular}
\end{table}

Invalid or non-finite raw candidates are excluded from selection.  They are not
silently repaired by a pseudoinverse, ridge penalty, or substitute method.  A
complete method failure would remain in the diagnostics and create missing
predictions rather than a dropped fold.  In the realised study, all 140 stored
fits and all predictions are finite.  These safeguards make clear why selected
component counts are method-specific indices rather than a common complexity
scale.

\section{Development evidence}
\label{sec:app-spectroscopy-development}

Table~\ref{tab:app-spectroscopy-development-rmse} gives the repeat-level
development results.  Arnoldi has the smallest mean repeated-CV RMSE, while
Raw has the second-smallest mean but a much larger standard deviation.  Raw's
five repeat RMSEs are 4.850, 4.960, 5.997, 4.844, and 4.760; the single 5.997
repeat creates its adverse tail.  The contiguous-fold check reverses the top
pair, with Raw at 4.913 and Arnoldi at 5.120.  Hence development prediction
supports Raw and Arnoldi as the leading pair but does not establish a stable
ordering between them.

\begin{table}[htbp]
    \centering
    \scriptsize
    \caption[Repeated-development prediction errors]{Development RMSE.  The first three columns summarise five complete
    repeated nested-CV out-of-fold estimates.  The final two columns are single
    prespecified sensitivities.}
    \label{tab:app-spectroscopy-development-rmse}
    \begin{tabular}{@{}lrrrrr@{}}
        \toprule
        Method & Mean (SD) & Median & Range & Contiguous folds & Original validation \\
        \midrule
        CG--FPLS     & 5.924 (0.083) & 5.940 & 5.784--5.990 & 6.149 & 5.998 \\
        Raw FPLS     & 5.082 (0.517) & 4.850 & 4.760--5.997 & 4.913 & 4.132 \\
        Arnoldi FPLS & 4.917 (0.131) & 4.954 & 4.741--5.084 & 5.120 & 4.150 \\
        FPCR         & 5.364 (0.219) & 5.396 & 5.074--5.624 & 5.349 & 3.978 \\
        \bottomrule
    \end{tabular}
\end{table}

The original 155-to-40 split gives the metrics in Table
\ref{tab:app-spectroscopy-original-split}.  All biases are positive, but the
small validation set and its restricted assay support make this a weak basis
for general calibration claims.  In particular, every $R^2$ is negative
because its SSE is compared with the unusually small validation SST.  This is
compatible with RMSE around four for three methods and does not contradict the
broader benchmark signal.

\begin{table}[htbp]
    \centering
    \small
    \caption[Original calibration--validation split sensitivity]{Original supplied-split sensitivity: fitting on the 155
    calibration tablets and evaluating on the 40 validation tablets.  Bias is
    prediction minus observation.}
    \label{tab:app-spectroscopy-original-split}
    \begin{tabular}{@{}lrrrr@{}}
        \toprule
        Method & RMSE & MAE & Bias & $R^2$ \\
        \midrule
        CG--FPLS     & 5.998 & 4.820 & 2.109 & -5.442 \\
        Raw FPLS     & 4.132 & 3.566 & 1.460 & -2.057 \\
        Arnoldi FPLS & 4.150 & 3.558 & 1.406 & -2.085 \\
        FPCR         & 3.978 & 3.431 & 1.512 & -1.834 \\
        \bottomrule
    \end{tabular}
\end{table}

Selection is more revealing than the small differences among the leading
development prediction errors.  Table
\ref{tab:app-spectroscopy-selection-stability} records every outer-fit range.
CG stops at three throughout.  Arnoldi remains tightly concentrated between
four and six.  Raw spans almost the entire requested path, reaches the cap in
two fits, and triggers the fixed instability rule in 24 of 25.  FPCR also
varies widely, although its spectral truncation index is not directly
comparable with a Krylov iteration.

\begin{table}[htbp]
    \centering
    \small
    \caption[Development selection and numerical stability]{Selection and numerical stability over 25 repeated-development
    outer fits.  Quartiles use the stored selected indices.}
    \label{tab:app-spectroscopy-selection-stability}
    \begin{tabular}{@{}lrrrrrr@{}}
        \toprule
        Method & Minimum & Q1 & Median & Q3 & Maximum & Cap / flags \\
        \midrule
        CG--FPLS     & 3 & 3  & 3  & 3  & 3  & 0 / 0 \\
        Raw FPLS     & 4 & 34 & 44 & 63 & 70 & 2 / 24 \\
        Arnoldi FPLS & 4 & 5  & 5  & 6  & 6  & 0 / 0 \\
        FPCR         & 7 & 12 & 21 & 39 & 57 & 0 / n/a \\
        \bottomrule
    \end{tabular}
\end{table}

\section{Benchmark uncertainty and residuals}
\label{sec:app-spectroscopy-uncertainty}

Chapter~\ref{chap:spectroscopy-application} gives all four primary point
estimates and marginal iid intervals.  Table
\ref{tab:app-spectroscopy-rmse-contrasts} supplies the complete prespecified
RMSE sensitivity for the three declared paired contrasts.  The focal
Raw-minus-Arnoldi interval remains below zero for every block length, as do the
larger CG-minus-Arnoldi and FPCR-minus-Arnoldi differences.  This order
robustness applies to conditional RMSE contrasts.  Bias conclusions are less
stable: the longer-block bias intervals for the secondary contrasts cross
zero.

\begin{table}[htbp]
    \centering
    \scriptsize
    \caption[Paired benchmark RMSE contrasts and intervals]{Declared benchmark RMSE contrasts and paired 95\% percentile
    intervals.  Circular blocks follow the supplied archive order.}
    \label{tab:app-spectroscopy-rmse-contrasts}
    \begin{tabular}{@{}lrrrrr@{}}
        \toprule
        Contrast & Estimate & Iid & Block 5 & Block 10 & Block 20 \\
        \midrule
        Raw $-$ Arnoldi
          & -0.0363 & [-0.0565, -0.0169] & [-0.0595, -0.0158]
          & [-0.0622, -0.0141] & [-0.0649, -0.0117] \\
        CG $-$ Arnoldi
          &  0.6500 & [ 0.3973,  0.9208] & [ 0.3286,  0.9783]
          & [ 0.2565,  1.0169] & [ 0.1664,  1.0687] \\
        FPCR $-$ Arnoldi
          &  0.5420 & [ 0.2855,  0.8442] & [ 0.2403,  0.8839]
          & [ 0.2015,  0.9218] & [ 0.1886,  0.9384] \\
        \bottomrule
    \end{tabular}
\end{table}

Figure~\ref{fig:app-spectroscopy-residuals} shows prediction-minus-observation
residuals against observed assay.  The methods share a few large errors,
including one high-assay tablet predicted near the centre of the range and a
low-assay tablet overpredicted by all procedures.  These rows remain in every
reported metric and in the bootstrap resampling population.  The common extremes widen marginal RMSE
intervals while largely cancelling in paired method differences.  Their
presence does not establish data error or a batch effect in the absence of
supporting metadata.

\begin{figure}[htbp]
    \centering
    \includegraphics[width=0.88\textwidth]{tablet_benchmark_residuals.pdf}
    \caption[Benchmark residuals against observed assay]{Designated-benchmark residuals, defined as prediction minus
    observation, plotted against observed assay.  All 460 tablets are retained.
    Shared large residuals are descriptive; the archive contains no batch or
    acquisition-order metadata with which to assign a physical cause.}
    \label{fig:app-spectroscopy-residuals}
\end{figure}

\section{Final-fit numerical diagnostics}
\label{sec:app-spectroscopy-numerics}

Table~\ref{tab:app-spectroscopy-final-diagnostics} contrasts final training fit
with numerical representation.  Raw and Arnoldi achieve nearly the same
training RMSE and coefficient norm, but through sharply different coordinate
systems.  Raw's selected condition numbers are at the reporting ceiling,
whereas Arnoldi retains an orthonormal basis to machine precision.  FPCR has
the smallest training RMSE and largest coefficient norm; these in-sample
quantities do not override the benchmark ordering.

\begin{table}[htbp]
    \centering
    \scriptsize
    \caption[Final-fit numerical diagnostics]{Final fit on all 195 development tablets.  Coefficient norms use
    the archived original-input vector under the $T^{-1}$ prediction
    convention.}
    \label{tab:app-spectroscopy-final-diagnostics}
    \begin{tabular}{@{}lrrrp{6.0cm}@{}}
        \toprule
        Method & $m$ & Training RMSE & $\|\widehat\beta\|_2$ & Principal diagnostic \\
        \midrule
        CG--FPLS & 3 & 5.544 & 18,997
          & Pilot $m=70$, $\widehat\sigma^2=10.467$; discrepancy threshold
            0.191 was crossed. \\
        Raw FPLS & 55 & 4.399 & 21,494
          & Basis and response-design conditions both recorded at
            $4.50\times10^{15}$; instability flag set. \\
        Arnoldi FPLS & 5 & 4.395 & 21,573
          & Basis condition 1; response-design condition 11.95; orthogonality
            defect $7.40\times10^{-16}$. \\
        FPCR & 53 & 3.400 & 45,334
          & Retained eigenvalue $7.02\times10^{-7}$; retained spectral
            condition 35,919. \\
        \bottomrule
    \end{tabular}
\end{table}

Figure~\ref{fig:app-spectroscopy-krylov-numerics} gives the complete relevant
paths.  For Raw, the basis condition is already $5.74\times10^8$ at $m=6$,
above the fixed $1/\sqrt{\epsilon_{\mathrm{mach}}}=6.71\times10^7$ threshold;
the response design is still worse.  The raw Krylov column norm falls by more
than 80 orders of magnitude by the selected $m=55$, and the aggregate inner-CV
criterion contains invalid candidate dimensions.  Its finite selected
prediction must therefore not be equated with stable continuation of the
intended power basis.

Arnoldi's basis condition remains one along the requested path, and its
orthogonality defect stays near machine precision.  The defect is exactly zero
at $m=1$, so that point is absent on the logarithmic vertical scale.  At the
selected $m=5$, the response-space least-squares problem is also well
conditioned.  Raw and Arnoldi nevertheless make almost identical benchmark
predictions: correlation 0.999934, RMS paired difference 0.174, and maximum
absolute difference 0.445 assay units.  The empirical role of orthogonalisation
is thus reliable representation of essentially the same predictive relation,
not guaranteed improvement of every realised loss.

\begin{figure}[htbp]
    \centering
    \includegraphics[width=0.82\textwidth]{tablet_raw_arnoldi_numerical_diagnostics.pdf}
    \caption[Final raw and Arnoldi Krylov diagnostics]{Final raw and Arnoldi Krylov diagnostics.  The upper panel shows
    raw-basis and response-design condition numbers against the fixed
    $1/\sqrt{\epsilon_{\mathrm{mach}}}$ flag threshold.  The lower panel shows
    Arnoldi's spectral-norm orthogonality defect; its exactly zero $m=1$ value
    is not displayed on the logarithmic scale.  Vertical markers identify the
    selected dimensions 55 and 5.}
    \label{fig:app-spectroscopy-krylov-numerics}
\end{figure}

\section{Preprocessing and instrument sensitivities}
\label{sec:app-spectroscopy-sensitivity}

Table~\ref{tab:app-spectroscopy-sensitivities} and Figure
\ref{fig:app-spectroscopy-sensitivity} report every prespecified benchmark
variant.  The full instrument-1 analysis remains the sole primary result.
Trimming the high-wavelength tail improves CG--FPLS, Arnoldi, and FPCR and
compresses the four-method RMSE range, but slightly worsens Raw.  Smoothing
changes both errors and selected indices without uniformly improving them.
Instrument~2 raises every RMSE; because it measures the same tablets, this is
measurement-system sensitivity rather than independent replication.

\begin{table}[htbp]
    \centering
    \scriptsize
    \caption[Complete prespecified benchmark sensitivities]{Complete prespecified benchmark sensitivities.  Parenthesised
    values in the RMSE column are selected component or stopping indices.
    Sensitivity rows are descriptive point estimates without variant-specific
    uncertainty intervals.}
    \label{tab:app-spectroscopy-sensitivities}
    \begin{tabular}{@{}llrrrr@{}}
        \toprule
        Variant & Method & RMSE ($m$) & MAE & Bias & $R^2$ \\
        \midrule
        Instrument 1, full
          & CG--FPLS     & 5.293 (3)  & 3.673 &  0.177 & 0.887 \\
          & Raw FPLS     & 4.607 (55) & 3.036 & -0.303 & 0.914 \\
          & Arnoldi FPLS & 4.643 (5)  & 3.086 & -0.301 & 0.913 \\
          & FPCR         & 5.185 (53) & 3.547 & -0.628 & 0.892 \\
        \addlinespace
        Instrument 1, 600--1798~nm
          & CG--FPLS     & 4.669 (3)  & 3.007 & -0.456 & 0.912 \\
          & Raw FPLS     & 4.767 (10) & 3.105 & -0.476 & 0.908 \\
          & Arnoldi FPLS & 4.632 (8)  & 2.846 & -0.513 & 0.913 \\
          & FPCR         & 4.719 (14) & 2.877 & -0.575 & 0.910 \\
        \addlinespace
        Instrument 1, SG smooth
          & CG--FPLS     & 5.341 (3)  & 3.717 &  0.013 & 0.885 \\
          & Raw FPLS     & 4.663 (18) & 3.056 & -0.497 & 0.912 \\
          & Arnoldi FPLS & 4.953 (9)  & 3.042 & -0.491 & 0.901 \\
          & FPCR         & 4.823 (27) & 3.053 & -0.444 & 0.906 \\
        \addlinespace
        Instrument 2, full
          & CG--FPLS     & 6.764 (3)  & 5.100 &  1.022 & 0.815 \\
          & Raw FPLS     & 5.042 (53) & 3.242 & -0.601 & 0.897 \\
          & Arnoldi FPLS & 5.095 (5)  & 3.279 & -0.562 & 0.895 \\
          & FPCR         & 6.395 (37) & 4.150 & -0.921 & 0.835 \\
        \bottomrule
    \end{tabular}
\end{table}

\begin{figure}[htbp]
    \centering
    \includegraphics[width=0.86\textwidth]{tablet_benchmark_sensitivity_rmse.pdf}
    \caption[Benchmark RMSE across prespecified sensitivities]{RMSE across the primary analysis and three prespecified
    sensitivities.  Points are descriptive benchmark estimates without
    uncertainty intervals.  Horizontal comparisons change the spectral
    representation or instrument and do not define a post hoc replacement for
    the primary full-spectrum instrument-1 result.}
    \label{fig:app-spectroscopy-sensitivity}
\end{figure}

The Raw--Arnoldi point ordering changes across variants, reinforcing that the
primary 0.036-unit difference is not a universal ranking.  Trimming improves
all methods except Raw, but this association does not prove that visible
high-wavelength spikes caused the primary residuals.  Similarly, poorer
instrument-2 results do not identify a transfer mechanism without a dedicated
calibration-transfer design.